\chapter{2022:Hugo Duminil-Copin}\label{ch2022hugoduminilcopin:hugo-duminil-copin-fields-medal-2022-a-comprehensive-report-on-his-work-and-impact}

\begin{center}\rule{0.5\linewidth}{0.5pt}\end{center}

\textbf{Prepared:} July 2026

\textbf{Subject:} Hugo Duminil-Copin (b.\ August 26, 1985, Essonne, France)

\textbf{Position:} Directeur de recherche, Institut des Hautes 'Etudes Scientifiques (IHES); Professor, University of Geneva

\textbf{Doctoral Advisor:} Stanislav Smirnov (University of Geneva)

\textbf{Fields Medal Citation:} \emph{``For solving longstanding problems in the critical theory of statistical mechanics, in particular the critical behaviour of the Ising model in three and more dimensions.''}

\begin{center}\rule{0.5\linewidth}{0.5pt}\end{center}

\begin{prereqs}
Probability theory (percolation, Markov chains, random walks), statistical mechanics (Ising model, Potts model, phase transitions), basic functional analysis, Fourier analysis.
\end{prereqs}

\begin{objectives}
Explain the critical exponents and phase transitions of the Ising model; describe sharp phase transition theorems for the Ising and Potts models; understand the relationship between the Ising model and the Potts model via the random-cluster representation.
\end{objectives}

\section{Problem History}

The Ising model, proposed by Wilhelm Lenz in 1920 and solved by his student Ernst Ising in 1925, is the simplest model of a ferromagnet. Spins take values in $\{+1, -1\}$ and sit on the vertices of a lattice, interacting with their neighbors. The partition function is $Z_{\beta}(h) = \sum_{\sigma} e^{\beta \sum_{xy \in E} \sigma_x \sigma_y + h \sum_{x \in V} \sigma_x}$, where $\beta > 0$ is the inverse temperature and $h \ge 0$ is the external field. This deceptively simple model captures the essential physics of magnetic phase transitions and has become one of the most studied models in mathematical physics.

The key parameter is the critical inverse temperature $\beta_c$. For $\beta < \beta_c$, the system is in the disordered phase with no spontaneous magnetization. For $\beta > \beta_c$, the system exhibits spontaneous symmetry breaking: $m^*(\beta) = \lim_{h \downarrow 0} \langle \sigma_0 \rangle_{\beta, h} > 0$. In 2 dimensions, Onsager (1944) solved the model exactly and showed the phase transition is continuous (second-order). The exact solution was a landmark achievement in mathematical physics and remains one of the few exactly solvable models in statistical mechanics.

In dimensions $d \ge 3$, the nature of the phase transition---whether it remains continuous---was one of the most important open problems in mathematical physics. The critical exponents $\beta, \gamma, \nu$ govern the behavior of physical quantities near the critical point: $m^* \sim (\beta - \beta_c)^\beta$, susceptibility $\chi \sim |\beta_c - \beta|^{-\gamma}$, correlation length $\xi \sim |\beta_c - \beta|^{-\nu}$. The universality hypothesis predicts these exponents depend only on dimension $d$, not on the details of the lattice or interaction.

The \textbf{Potts model}, a generalization with $q \ge 2$ spin values, has a first-order phase transition for $q > 4$ in 2D (Baxter). Proving the sharpness of the phase transition for the Potts model in higher dimensions was a major open problem. The Potts model interpolates between the Ising model ($q = 2$) and the percolation model ($q \to 1$), providing a unifying framework for understanding different types of phase transitions.

The \textbf{random-cluster representation} (Fortuin--Kasteleyn, 1972) provides a powerful tool for studying the Potts model. The partition function can be rewritten as $Z_{\text{Potts}} = \sum_{A \subseteq E} p^{|A|}(1-p)^{|E|-|A|} q^{k(A)}$, where $k(A)$ is the number of connected components in the subgraph $(V, A)$. This representation connects the Potts model to percolation theory and allows the use of tools from probability theory.

The \textbf{correlation length} $\xi(\beta)$ measures the exponential decay rate of the two-point function. At the critical point $\beta = \beta_c$, the correlation length diverges: $\xi(\beta) \to \infty$ as $\beta \to \beta_c$. The rate of divergence is governed by the critical exponent $\nu$: $\xi(\beta) \sim |\beta - \beta_c|^{-\nu}$. Numerical simulations suggest $\nu \approx 0.6300$ for the 3D Ising model, but a rigorous proof of any specific value was unknown.

The \textbf{Lee--Yang theorem} (1952) implies that all zeros of the partition function as a function of the fugacity lie on the unit circle, constraining the possible types of phase transitions. The theorem has been generalized to other spin systems and provides a powerful tool for studying phase transitions.

\section{How Was It Solved}

\textbf{The Ising model has no first-order transition (2016).} With Kozma and Schramm, Duminil-Copin proved that the 3D Ising model has no first-order phase transition: the magnetization curve is continuous. The proof uses a new version of the Russo--Seymour--Welsh (RSW) theorem adapted to the Ising setting, together with random-cluster representations. The key idea is to show that crossing probabilities at criticality are bounded away from $0$ and $1$, which implies the absence of a first-order transition.

The RSW theorem is a cornerstone of 2D percolation theory, providing sharp bounds on crossing probabilities in rectangles at criticality. Extending it to the Ising model required developing new techniques because the Ising model has long-range correlations that are absent in percolation. The key insight was to use the random-cluster representation, which expresses the Ising model as a correlated percolation model, and then to develop new comparison inequalities that relate crossing probabilities in the random-cluster model to those in standard percolation.

The absence of a first-order transition has profound physical implications: it means that the magnetization changes continuously as the temperature passes through the critical point. This is consistent with the physical picture of a second-order phase transition, where the system undergoes a continuous change in its macroscopic properties. The result also provides a starting point for the study of critical exponents and universality classes in 3D statistical mechanics.

The RSW theorem, originally proved for percolation by Russo (1978) and Seymour--Welsh (1978), provides bounds on crossing probabilities in rectangles. Duminil-Copin and collaborators adapted this theorem to the Ising model by using the random-cluster representation and the FKG inequality. The key technical innovation is a \textbf{decoupling inequality} that relates crossing probabilities at different scales.

\textbf{Sharp threshold for Potts models (2017).} With Tassion, a sharp threshold result was established for the survival probability of random walks in random environments. The proof relies on comparing the Potts model to a random-cluster model and analyzing crossing probabilities using RSW-type estimates.

\textbf{Hypercontractivity and correlation inequalities.} Duminil-Copin and Gorban proved a new inequality relating the two-point function to the spectral gap using hypercontractivity of the semigroup generated by the Ising Glauber dynamics. The Glauber dynamics is a Markov chain that updates one spin at a time. The spectral gap controls the rate of convergence to equilibrium: the mixing time is $O(\log n / \gamma)$, where $\gamma$ is the spectral gap. The hypercontractivity inequality provides a sharp connection between the dynamics and the statics of the model.

The hypercontractivity inequality for the Ising model states that for any function $f$ on the spin configurations,
\[
\|P_t f\|_{q} \le \|f\|_{p}
\]
where $P_t$ is the semigroup generated by the Glauber dynamics and the norms are taken with respect to the Gibbs measure. The key parameter is the relationship between $p$, $q$, and $t$: the inequality holds for $q \ge p \ge 1$ when $t$ is large enough, with the optimal time depending on the spectral gap. This inequality implies that the correlations between spins at different sites decay exponentially in the distance, which is the hallmark of the disordered phase.

The hypercontractivity technique has been applied to other spin systems, including the Potts model and continuous-spin models. The key advantage is that it provides sharp bounds on correlation functions that are difficult to obtain by other methods. The technique has also been used to study the mixing time of Markov chains on graphs, with applications to computer science and statistical mechanics.

\textbf{The sharp critical point (2019).} With Raoufi and Tassion, Duminil-Copin proved that the critical point of the Ising model on $\mathbb{Z}^d$ equals $\beta_c = \operatorname{atanh}(1/(2d-1))$, resolving a conjecture of Coniglio and Klein. The proof uses a self-duality argument: the model on the dual lattice satisfies $\sinh(2\beta) \sinh(2\beta^*) = 1$, and at criticality $\beta = \beta^*$.

The critical inverse temperature $\beta_c$ is the threshold between the disordered phase ($\beta < \beta_c$) and the ordered phase ($\beta > \beta_c$). The exact value $\beta_c = \operatorname{atanh}(1/(2d-1))$ was conjectured by Coniglio and Klein based on physical arguments involving the duality transformation. The proof uses the random-cluster representation and the FKG inequality to establish sharp bounds on the two-point function at the conjectured critical point. The key estimate is that the two-point function decays polynomially at $\beta_c$, which implies that the correlation length diverges and the system is at criticality.

\textbf{The lace expansion for the Potts model (2022).} With Aizenman, Duminil-Copin, and Harel applied the lace expansion to the $q$-state Potts model above the upper critical dimension $d_c = 4$. The proof shows that the two-point function decays like a simple random walk, implying Gaussian critical exponents.

\section{Impact of This Discovery}

The continuity of the phase transition for the 3D Ising model was one of the most important open problems in mathematical physics. Duminil-Copin's work provides a complete picture of the Ising model's critical behavior, together with Smirnov's earlier 2D result. The results have been celebrated as a major breakthrough in mathematical physics.

The hypercontractivity techniques have been applied to other spin systems, including the clock model and continuous-spin models. The sharp threshold methods have found applications in percolation theory and the analysis of random graphs. The RSW theorem for the Ising model has been generalized to other models.

Duminil-Copin's work bridges probability theory and mathematical physics, providing rigorous tools for studying critical phenomena that physicists had long conjectured but not proven. The connection between the Ising model and conformal field theory in dimension 2, established by Smirnov, has been extended to a more rigorous framework.

The \textbf{Schramm--Loewner evolution} (SLE) describes the scaling limit of interfaces in 2D statistical mechanics models. Smirnov's proof that the percolation interface converges to SLE$_6$ was a landmark result. Duminil-Copin's work provides tools for extending this connection to other models in 2D.

The \textbf{random-cluster representation} provides a powerful framework for studying the Potts model. The partition function of the $q$-state Potts model can be rewritten as
\[
Z_{\text{Potts}} = \sum_{A \subseteq E} p^{|A|}(1-p)^{|E|-|A|} q^{k(A)},
\]
where $p = 1 - e^{-\beta}$ is the edge probability, $\beta$ is the inverse temperature, and $k(A)$ is the number of connected components in the subgraph $(V, A)$. This representation connects the Potts model to percolation theory and allows the use of tools from probability theory. The random-cluster model has the important property that it satisfies the FKG inequality, which implies that increasing events are positively correlated. This property is essential for the proof of the RSW theorem and for the analysis of crossing probabilities at criticality.

The \textbf{Lee--Yang theorem} (1952) states that all zeros of the partition function of the Ising model as a function of the fugacity $z = e^{h}$ lie on the unit circle $|z|=1$. This theorem constrains the possible types of phase transitions: if the zeros accumulate at $z=1$, the system has a continuous phase transition; if they are separated from $z=1$, the system is in the disordered phase. The Lee--Yang theorem has been generalized to other spin systems and provides a powerful tool for studying phase transitions.

\section{Modern Views and Improvements}

Duminil-Copin continues to develop the mathematical framework for understanding phase transitions. Recent work with Klenke on the Ising model on the hyperbolic plane explores the interplay between geometry and phase transitions: there is a first-order transition at finite temperature for the Ising model on the hyperbolic plane.

Connections to conformal field theory in dimension 2 are well-established. Extending these to dimension 3 remains a major open challenge. The conformal bootstrap, a non-perturbative method from physics that predicts critical exponents, has been verified numerically in 3D but lacks rigorous mathematical foundations.

The Ising model in random environments, long-range interactions, and on non-amenable groups represent current frontiers. The lace expansion has been extended to the $\varphi^4$ lattice field theory, providing a rigorous renormalization group analysis of the model.

The \textbf{random-cluster model} with general $q$ continues to be an active area. The critical exponents for $q > 1$ are not known rigorously in 3D. The connection between the random-cluster model and the $O(n)$ loop model provides a unifying framework for studying phase transitions.

The \textbf{continuum limit} of the 3D Ising model is a central open problem. In 2D, the continuum limit is described by a conformal field theory with central charge $c = 1/2$. In 3D, no rigorous construction exists. Duminil-Copin's work provides the discrete-level foundations needed for any rigorous construction.

Current research includes the study of the Ising model on general graphs, the development of new correlation inequalities, and the extension of RSW techniques to non-planar settings. The interplay between Duminil-Copin's probabilistic methods and the field-theoretic approach to critical phenomena continues to produce new results.

The Ising model on non-amenable graphs (such as hyperbolic planes and trees) exhibits qualitatively different behavior from the model on $\mathbb{Z}^d$. There is a first-order phase transition at finite temperature, and the critical exponents are different from those in Euclidean space. Duminil-Copin's work with Klenke on the Ising model on the hyperbolic plane has provided the first rigorous results on this problem, opening new directions in the study of phase transitions on non-Euclidean spaces.

The lace expansion, originally developed for percolation and the Ising model above the upper critical dimension, has been extended to a wider class of models. Duminil-Copin and collaborators have applied the lace expansion to the $O(n)$ loop model and the $\varphi^4$ lattice field theory, obtaining sharp results on critical behavior in high dimensions. These results provide a rigorous foundation for the renormalization group theory of critical phenomena.

The study of the Ising model on general graphs, including random graphs and graphs with long-range interactions, represents another active frontier. The critical behavior on random graphs such as the Erd\"{o}s--R\'{e}nyi model and scale-free networks differs qualitatively from the behavior on Euclidean lattices, and Duminil-Copin's techniques have been adapted to study phase transitions in these settings.

The connection between the Ising model and conformal field theory in dimension 2, established by Smirnov's proof of conformal invariance at criticality, has been further developed by Duminil-Copin's work. In dimension 3, the conformal bootstrap predicts the values of critical exponents, but a rigorous mathematical foundation for these predictions is still lacking. The development of rigorous methods for studying the conformal properties of the 3D Ising model is an active area that builds on Duminil-Copin's results on the continuity of the phase transition and the sharp determination of the critical point.

The study of the Potts model in dimensions greater than 4, above the upper critical dimension, has been influenced by Duminil-Copin's work. The lace expansion, which provides sharp results on critical behavior above the upper critical dimension, has been applied to the Potts model and to other spin systems. The development of new techniques for applying the lace expansion to models with complex phase structures is an active area that builds on Duminil-Copin's foundational work on the critical behavior of the Ising model.

\begin{center}\rule{0.5\linewidth}{0.5pt}\end{center}

\section*{Summary}
\addcontentsline{toc}{section}{Summary}

Duminil-Copin proved the continuity of the phase transition for the 3D Ising model using random-cluster representations, RSW theorems, and hypercontractivity techniques. His work provides the rigorous mathematical foundation for understanding critical phenomena.
